\section{Nguy cơ đối với tính hợp lệ}
\label{sec:threats-limitations}

\paragraph{Tính hợp lệ nội tại.}
Prompt, mô hình và judge đều có thể ảnh hưởng tới điểm số. Chúng tôi giảm rủi ro này
bằng cách lưu manifest cho từng lượt chạy, giữ cấu hình decoding giống nhau giữa các
tác tử và kích thước, tách mô hình sinh khỏi mô hình judge và sử dụng hai judge. Một
yếu tố gây nhiễu vẫn còn: AlignmentGraph-ISD dùng nhiều lời gọi mô hình hơn Baseline
một-prompt, nên một phần lợi thế trên điểm judge có thể đến từ test-time compute bổ
sung. Nghiên cứu báo số lời gọi và token theo tác tử trong artifact
(\TBD{bảng resource theo tác tử}), nhưng chưa có đối sánh cân bằng ngân sách; kết luận
vì vậy áp dụng cho các tác tử như chúng được thiết kế, không phải cho các phương án đã
cân bằng compute.

\paragraph{Tính hợp lệ khái niệm.}
Rubric ADDIE và điểm trajectory đo chất lượng đầu ra cùng dấu vết thực thi, không đo
kết quả học tập. Thước đo alignment của RQ2 khắc phục việc alignment không được đo
trực tiếp, nhưng mang các giới hạn riêng cần khai báo. Thứ nhất, bộ phân loại Bloom và
encoder có sai số ($\kappa<1$ so với nhãn chuyên gia trong các nghiên cứu
gốc~\cite{waheed2021bloomnet,li2022bloomclassification}); vì cùng một pipeline được áp
dụng cho mọi tác tử, sai số này chủ yếu ảnh hưởng giá trị tuyệt đối, nên kết quả RQ2
được diễn giải là \emph{xếp hạng tương đối giữa các tác tử}, không phải một cuộc kiểm
toán alignment tuyệt đối. Thứ hai, các tài nguyên phân loại Bloom chủ yếu được huấn
luyện trên tiếng Anh; việc chuyển sang scenario tiếng Hàn (qua encoder đa ngữ và
fallback từ điển tiếng Hàn) là một nguy cơ chuyển giao chưa được kiểm chứng độc lập.
Thứ ba, thước đo lượng hoá alignment cấu trúc--ngữ nghĩa giữa các thành phần văn bản;
nó không đo tính đúng đắn sư phạm sâu hơn hay kết quả học tập thực tế. Một đánh giá
với chuyên gia gán nhãn vẫn là bước xác nhận cần thiết.

\paragraph{Tính hợp lệ ngoại tại.}
Các scenario đều là dữ liệu tổng hợp từ một benchmark; kết quả có thể thay đổi trên
brief thực tế, lĩnh vực khác hoặc benchmark ISD khác. Dải kích thước nằm trong một họ
mô hình duy nhất (Qwen) trên cùng runtime: thiết kế này cô lập hiệu ứng kích thước
khỏi khác biệt family, nhưng đổi lại, kết luận về model scale chỉ chắc chắn trong họ
mô hình đã đánh giá và cần được kiểm tra lại trên các họ khác. Ngoài ra, dù động cơ
hướng tới tự lưu trữ, các thí nghiệm gọi mô hình qua API bên ngoài; hành vi trên
hạ tầng on-premise với quantization thực tế chưa được đo.

\paragraph{Tính hợp lệ của kết luận thống kê.}
Các hiệu ứng được ước lượng trên tối đa 90 cặp scenario với ba lượt chạy mỗi tổ hợp;
SD giữa các lượt được báo cáo và các kiểm định dùng trung bình theo scenario, nên
run-to-run variability được lượng hoá nhưng chưa được mô hình hoá đầy đủ (ví dụ bằng
mixed model theo cả scenario và lượt chạy). Complete-case có thể tạo thiên lệch nếu
lỗi không xảy ra ngẫu nhiên; phân tích độ nhạy với scenario lỗi bị phạt điểm
(\cref{sec:stats}) kiểm tra mức nhạy của kết luận với chính sách này. Các scenario
trong cùng lĩnh vực có thể phụ thuộc lẫn nhau; bootstrap theo cụm lĩnh vực là một
kiểm tra bổ sung đáng thực hiện.

\section{Hạn chế và hướng phát triển}
\label{sec:future-work}

Nghiên cứu giới hạn ở việc sinh một chương trình từ một design brief cho trước. Brief
trong thực tế thường hình thành từ trao đổi với các bên liên quan, nên một giao diện hội
thoại để thu thập brief là mở rộng tự nhiên. Harness cũng có thể tiếp nhận tri thức tổ
chức---chẳng hạn rubric nội bộ hoặc quy ước chuyên ngành---qua retrieval, nhưng khả năng
này chưa được đánh giá. Về đo lường, ưu tiên tiếp theo là hiệu chuẩn thước đo alignment
với nhãn của chuyên gia ISD (báo agreement giữa người gán nhãn) và liên hệ các chỉ số
alignment với kết quả học tập thực tế. Về triển khai, một đánh giá trên hạ tầng
on-premise và một nghiên cứu người dùng với chuyên gia thiết kế giảng dạy sử dụng
artifact đồ thị để audit là hai bước cần thiết để kiểm chứng các lợi ích đã nêu trong
động cơ.
